\documentclass[11pt,a4paper]{article}
\usepackage[utf8]{inputenc}
\usepackage[T1]{fontenc}
\usepackage[margin=2.2cm]{geometry}
\usepackage{tabularx}
\usepackage{booktabs}
\usepackage{xcolor}
\usepackage{colortbl}
\usepackage{hyperref}
\usepackage{enumitem}
\usepackage{longtable}
\hypersetup{colorlinks=true,linkcolor=black,urlcolor=blue!60!black}

\definecolor{hdr}{RGB}{28,42,64}
\definecolor{rowgray}{RGB}{243,245,248}
\definecolor{good}{RGB}{20,110,60}
\definecolor{bad}{RGB}{160,35,35}

\newcolumntype{Y}{>{\raggedright\arraybackslash}X}
\newcommand{\yes}{\textcolor{good}{\textbf{yes}}}
\newcommand{\no}{\textcolor{bad}{\textbf{no}}}

\title{\vspace{-1.5cm}\textbf{Choosing a Fine-Tuning Method for a Graph-RAG\\Enterprise Assistant}\\[0.3cm]
\large A measured comparison against this system's own evaluation data}
\author{Prepared for Iyed Mediouni --- PFE, TEK-UP $\times$ TELNET Holding}
\date{15 August 2026}

\begin{document}
\maketitle
\vspace{-0.8cm}

\section*{Purpose}
This report decides which fine-tuning method --- if any --- is justified for the assistant, using
evidence measured on this system rather than general claims. It scores every candidate against
constraints that emerged during development, including the supervising professor's position that
code fixes remove the need for fine-tuning.

\section{A distinction that changes the question}

\textbf{RAFT is not an alternative to QLoRA.} They sit on different axes:

\begin{itemize}[leftmargin=1.2em,itemsep=2pt]
  \item \textbf{Training objective / data recipe} --- what the model is taught:
        SFT, RAFT, DPO, ORPO, KTO, SimPO, GRPO, RLHF.
  \item \textbf{Parameter-efficiency method} --- how the weights are updated:
        full fine-tuning, LoRA, QLoRA, DoRA, GaLore.
\end{itemize}

RAFT is run \emph{with} QLoRA. The choice is therefore not ``RAFT or QLoRA'' but
``which objective, executed with which parameter method''. This matters for the defence: a jury
member who knows the field will notice the conflation immediately.

\section{Timeline of the techniques considered}

\begin{table}[h]\centering\small
\begin{tabularx}{\textwidth}{l l Y}
\toprule
\textbf{Year} & \textbf{Technique} & \textbf{Contribution} \\
\midrule
2021 & LoRA & Low-rank adapters; freezes the base model (Microsoft) \\
\rowcolor{rowgray} 2022 & RLHF / PPO & Preference alignment via a learned reward model (InstructGPT) \\
2023 & QLoRA & 4-bit NF4 base + LoRA; 65B-class tuning on one GPU (U.~Washington) \\
\rowcolor{rowgray} 2023 & DPO & Preference tuning without a reward model (Stanford) \\
2024 & RAFT & SFT with distractor documents + cited chain-of-thought (UC Berkeley/Microsoft) \\
\rowcolor{rowgray} 2024 & DoRA & Splits updates into magnitude and direction; $+1$--$4.4\%$ over LoRA (NVIDIA) \\
2024 & KTO / ORPO / SimPO & Preference variants; ORPO merges SFT and DPO into one stage \\
\rowcolor{rowgray} 2024--25 & GRPO & Group-relative RL for verifiable rewards; popularised by DeepSeek-R1 \\
2025 & GraphRAFT & RAFT for knowledge graphs; trains correct graph-query generation \\
\rowcolor{rowgray} 2025 & DTA (Divide-Then-Align) & Four knowledge quadrants; teaches calibrated ``I don't know'' (ACL 2025) \\
2025 & CRAFT / ALoFTRAG / RbFT & Compute-efficient, local-only and noise-robust RAFT variants \\
\rowcolor{rowgray} 2026 & TRL v1.0 & One library unifying SFT/DPO/KTO/ORPO/GRPO trainers (April 2026) \\
\bottomrule
\end{tabularx}
\end{table}

\section{What this system actually measured}

Every score below is anchored on these figures, obtained on this assistant's own corpus.

\begin{table}[h]\centering\small
\begin{tabularx}{\textwidth}{Y r}
\toprule
\textbf{Observation} & \textbf{Value} \\
\midrule
Judged benchmark records (RAGAS + validator, no approximations) & 299 \\
\rowcolor{rowgray} Faithfulness $<0.70$ & 13.7\,\% $\pm$ 3.9 \\
Quality grades GOLD / SILVER / BRONZE / FAILED & 229 / 51 / 16 / 3 \\
\rowcolor{rowgray} Fabrication defects (deterministic validator, checkable answers) & 10.0\,\% local, 7.1\,\% fp16 \\
Share of fabrication defects that are invented URLs & $\approx$95\,\% \\
\rowcolor{rowgray} Overlap between validator-detected and judge-detected defects & 3 of 63 \\
Refusals that are ``bare'' (decline, offer nothing usable) & 41\,\% -- 57\,\% \\
\rowcolor{rowgray} Adapter~1 (DPO, 188 pairs): gate verdict & \textbf{REJECTED} \\
\quad source coverage 11.5\,\% $\rightarrow$ 16.4\,\%, refusals 75.4\,\% $\rightarrow$ 93.4\,\% (rule D8 tripped) & \\
\rowcolor{rowgray} Natural preference pairs available (real text on both sides) & 96 \\
On-policy SFT demonstrations available & 419 \\
Retrieval demonstrably returns distractors (asked 4.16.0; returned 4.16.1, 4.28.0) & yes \\
\bottomrule
\end{tabularx}
\end{table}

\section{Scoring rubric}

Criteria are taken from constraints stated during this project, not invented for the report.
Each method scores 0--5 per criterion; the weight reflects how decisive that criterion is here.

\begin{table}[h]\centering\small
\begin{tabularx}{\textwidth}{c Y c}
\toprule
\textbf{\#} & \textbf{Criterion} & \textbf{Weight} \\
\midrule
C1 & \textbf{Not achievable by a code fix.} The professor's position: if a filter fixes it, weights should not. & $\times$3 \\
\rowcolor{rowgray} C2 & \textbf{Targets a defect measured in this system}, not a generic weakness. & $\times$3 \\
C3 & \textbf{Training data exists without synthesis.} A pool that is 93\,\% corrupted answers is weak evidence. & $\times$2 \\
\rowcolor{rowgray} C4 & \textbf{Improvement is detectable} on the available holdout with free, deterministic metrics. & $\times$2 \\
C5 & \textbf{Runs from the existing UI} without new training code. & $\times$2 \\
\rowcolor{rowgray} C6 & \textbf{Fits the compute budget} (6\,GB local; Kaggle T4$\times$2, 30\,h/week). & $\times$1 \\
C7 & \textbf{Defensible to a jury} --- published, citable, reproducible. & $\times$2 \\
\rowcolor{rowgray} C8 & \textbf{Matches the stated product preference}: a clean ``not available'' beats a tangential citation. & $\times$2 \\
\bottomrule
\end{tabularx}
\end{table}

\section{Scored comparison}

\begin{table}[h]\centering\footnotesize
\setlength{\tabcolsep}{4pt}
\begin{tabular}{l ccccccc c r}
\toprule
\textbf{Method (objective)} & \textbf{C1} & \textbf{C2} & \textbf{C3} & \textbf{C4} & \textbf{C5} & \textbf{C6} & \textbf{C7} & \textbf{C8} & \textbf{Total} \\
 & $\times$3 & $\times$3 & $\times$2 & $\times$2 & $\times$2 & $\times$1 & $\times$2 & $\times$2 & /85 \\
\midrule
Vanilla SFT (QLoRA) on GOLD answers & 2 & 1 & 4 & 2 & 5 & 5 & 3 & 2 & \textbf{47} \\
\rowcolor{rowgray} DPO on corruption pairs \emph{(built, rejected)} & 1 & 2 & 1 & 2 & 5 & 5 & 3 & 2 & \textbf{41} \\
DPO on the 96 natural pairs & 3 & 3 & 5 & 2 & 5 & 5 & 3 & 3 & \textbf{57} \\
\rowcolor{rowgray} ORPO / SimPO / KTO & 3 & 2 & 4 & 2 & 2 & 4 & 4 & 3 & \textbf{48} \\
GRPO (verifiable reward) & 4 & 2 & 5 & 3 & 1 & 1 & 4 & 3 & \textbf{50} \\
\rowcolor{rowgray} RLHF / PPO + reward model & 4 & 2 & 1 & 3 & 1 & 1 & 4 & 3 & \textbf{42} \\
\textbf{RAFT (SFT + distractors + cited CoT)} & 5 & 5 & 4 & 4 & 5 & 4 & 5 & 3 & \textbf{74} \\
\rowcolor{rowgray} \textbf{RAFT $\rightarrow$ DTA (two stage)} & 5 & 5 & 4 & 4 & 5 & 4 & 5 & 5 & \textbf{78} \\
GraphRAFT & 5 & 4 & 2 & 3 & 3 & 3 & 4 & 3 & \textbf{63} \\
\bottomrule
\end{tabular}
\end{table}

\noindent\textit{Parameter-efficiency axis:} QLoRA (2023) remains correct here --- 8B model, 6\,GB local
VRAM, T4 on Kaggle, and adapters that swap at serve time. DoRA (2024) offers $+1$--$4.4\%$ over LoRA
and is a drop-in upgrade, but it changes nothing about which objective to choose.

\section{Why the two leaders win}

\subsection*{RAFT (2024) --- scores 74}

RAFT trains the model on a question plus a \emph{golden} document and several \emph{distractor}
documents, with a chain-of-thought answer that quotes the supporting span verbatim. In a fraction
$(1-P)$ of examples the golden document is withheld entirely.

It scores highest on C1 and C2 because it targets a defect this system demonstrably has and no
filter can repair: \textbf{choosing correctly among retrieved documents}. When the assistant was
asked about ``Release notes 4.16.0'' and retrieval returned 4.16.1 and 4.28.0, the decision of what
to do with those near-misses is a model behaviour, not a string operation.

The published results also supply the strongest single argument against the ``just fix the code''
position: in the RAFT paper, \textbf{domain fine-tuning plus RAG is often worse than domain
fine-tuning alone} (HuggingFace 61.06 $\rightarrow$ 42.59; TensorFlow 86.56 $\rightarrow$ 60.29).
The interaction between tuning and retrieval is a genuine research problem, and RAFT is the
published remedy --- gains of up to 35 points over base+RAG on HotpotQA.

\subsection*{DTA / Divide-Then-Align (ACL 2025) --- lifts the total to 78}

RAFT has a documented weakness: it conditions models to answer even when the context is entirely
noise, which raises hallucination risk. DTA addresses exactly this. It divides samples into four
quadrants by whether the answer is available in the retrieved passages and in the model's internal
knowledge, then builds preference data per quadrant so the model learns to reply
\emph{``I don't know''} when the query falls outside both boundaries.

This is the only candidate that scores 5 on C8, because it optimises for the behaviour already
stated as preferred: a clean, short ``the data is not there'' rather than a citation of something
adjacent. It uses \textbf{DPO} --- so the method is not rejected, only redirected: the earlier
failure came from building preference pairs by corrupting answers, not from DPO itself.

\section{UI integration}

\begin{table}[h]\centering\small
\begin{tabularx}{\textwidth}{l c Y}
\toprule
\textbf{Method} & \textbf{New trainer?} & \textbf{What changes} \\
\midrule
RAFT & \no & Uses the existing \texttt{qlora} path; only the dataset builder differs \\
\rowcolor{rowgray} DTA & \no & Uses the existing \texttt{dpo} path; only the dataset builder differs \\
GRPO / RLHF & \yes & New trainer, reward plumbing, online rollouts \\
\rowcolor{rowgray} ORPO / SimPO / KTO & \yes & New TRL trainer classes to wire in \\
\bottomrule
\end{tabularx}
\end{table}

The training panel already exposes \texttt{method: [qlora, dpo, sequential]}, logs to MLflow, and can
activate an adapter as \texttt{FROM qwen3:8b + ADAPTER}. RAFT and DTA therefore require
\textbf{no new training code} --- they are data recipes that run through paths that already exist and
are already reachable from the interface. Every other candidate requires a new trainer.

\section{Recommendation}

\begin{enumerate}[leftmargin=1.4em,itemsep=3pt]
  \item \textbf{Adopt RAFT as the training objective, executed with QLoRA.} It is the only method
        scoring 5 on both ``no code fix can do this'' and ``targets a measured defect''.
  \item \textbf{Follow with a DTA stage on the existing \texttt{dpo} path} if time allows, to hold the
        abstention behaviour that is preferred for this product.
  \item \textbf{Reject GRPO, RLHF and a learned reward model} --- and say why: the reward signal here
        is already exact, deterministic and free (refusal classifier plus validator penalties). A
        learned reward model would add compute and approximation error where an exact reward exists.
  \item \textbf{Do the code fix as well.} Filtering links absent from the context is cheaper and more
        reliable than any weight update, and implementing it strengthens rather than weakens the
        argument for fine-tuning what remains.
\end{enumerate}

\section{Risks and honest limitations}

\begin{itemize}[leftmargin=1.2em,itemsep=2pt]
  \item RAFT's reported gains come from PubMed, HotpotQA and Gorilla --- not from this corpus.
        Transfer is plausible, not guaranteed.
  \item Generating cited chain-of-thought answers requires a stronger model for the training data
        (the paper used GPT-4-1106; the local equivalent here is the hosted judge), which reintroduces
        a quota dependency for dataset construction, though not for evaluation.
  \item The gate may still reject the result. That outcome remains a legitimate finding, and the
        first rejection is already documented with a diagnosed cause.
\end{itemize}

\section*{References}
\small
\begin{itemize}[leftmargin=1.2em,itemsep=1pt]
  \item Hu et al. \textit{LoRA: Low-Rank Adaptation of Large Language Models}, 2021.
  \item Ouyang et al. \textit{Training language models to follow instructions with human feedback} (InstructGPT), 2022.
  \item Dettmers et al. \textit{QLoRA: Efficient Finetuning of Quantized LLMs}, 2023.
  \item Rafailov et al. \textit{Direct Preference Optimization}, 2023.
  \item Zhang, Patil et al. \textit{RAFT: Adapting Language Model to Domain Specific RAG}, 2024. \url{https://arxiv.org/abs/2403.10131}
  \item Liu et al. \textit{DoRA: Weight-Decomposed Low-Rank Adaptation}, 2024.
  \item Shao et al. \textit{DeepSeekMath} (GRPO), 2024; popularised by DeepSeek-R1, 2025.
  \item \textit{Divide-Then-Align: Honest Alignment based on the Knowledge Boundary of RAG}, ACL 2025. \url{https://arxiv.org/pdf/2505.20871}
  \item Joren et al. \textit{Sufficient Context: A New Lens on Retrieval Augmented Generation}, 2024. \url{https://arxiv.org/pdf/2411.06037}
  \item \textit{Let Them Down Easy! Contextual Effects of LLM Guardrails}, EMNLP Findings 2025. \url{https://arxiv.org/abs/2506.00195}
\end{itemize}

\end{document}
